% Generated draft: requires mathematical and source review.
\documentclass[11pt]{article}
\usepackage[margin=25mm]{geometry}
\usepackage{fontspec}
\setmainfont{Latin Modern Roman}
\usepackage{amsmath,amssymb,mathtools,unicode-math}
\setmathfont{Latin Modern Math}
\usepackage{enumitem,needspace,xurl,hyperref}
\setlength{\parindent}{0pt}
\setlength{\parskip}{5pt}
\setlength{\emergencystretch}{3em}
\newcommand{\N}{\mathbb{N}}
\newcommand{\Z}{\mathbb{Z}}
\newcommand{\Q}{\mathbb{Q}}
\newcommand{\R}{\mathbb{R}}
\newcommand{\C}{\mathbb{C}}
\newcommand{\F}{\mathbb{F}}
\newcommand{\A}{\mathbb{A}}
\newcommand{\PP}{\mathbb{P}}
\newcommand{\E}{\mathbb{E}}
\newcommand{\eps}{\varepsilon}
\newcommand{\dd}{\,\mathrm{d}}
\begin{document}
\section*{UnsolvedMath ID: 9000002}
\textbf{Source problem code:} catalogue entry
\subsection{Definitions and mathematical statement}
The problem asks for a higher-dimensional analogue of the source paper's planar cropping and summation argument. In dimension three it suggests indexing an expected sum by quadruples $v_1,v_2,v_3,v_4$, with
\[
P=\operatorname{conv}\{0,v_1,v_2,v_3,v_4\},
\]
subject to an empty-lattice-point condition.

In the intended lattice setting one would take $v_i\in\Z^3$, but the literal phrase ``contains no lattice points'' cannot mean $P\cap\Z^3=\varnothing$: $0$ and the listed lattice points already belong to $P$. Possible intended conditions are absence of interior lattice points, $\operatorname{int}(P)\cap\Z^3=\varnothing$, or absence of any lattice points except the specified vertices. They differ, especially on boundary faces, and are not silently identified. The record does not say that all five listed points are distinct or vertices, or that $P$ is full dimensional. Five listed points in three-space also need not form a tetrahedron.

The crop operation, class of convex bodies, summand, normalization, multiplicities or symmetry quotient, summation order/convergence, and target identity are not given. Therefore no concrete series or constant is reconstructed from the paper title or the summary's schematic geometry-of-numbers expression. The proposed quadruple index set and the need to extend the planar argument are retained as the actual research direction, with the missing empty-polytope convention explicit.
\subsection{Short English statement}
Extend the planar crop-and-sum method to higher dimensions, especially a three-dimensional sum over specified lattice quadruples. Clarify the summand and whether empty means no interior lattice points or no extra lattice points at all.
\subsection{Sources}
\begin{itemize}
\item Local UnsolvedMath registry, permanent ID 9000002, exact source code catalogue entry, ``Higher-dimensional cropping formula''; attributed to Nikita Kalinin and Mikhail Shkolnikov, no proposed year supplied. Registry status open, classification OPEN-TRIAGE.
\item Kalinin and Shkolnikov, \emph{The Number pi and a Summation by SL(2,Z)} (2018), section 4, Questions, direction 2, \url{https://amj.math.stonybrook.edu/html-articles/Files-2015-2024/17-75/index.html}. Supplied HTML-text extraction accessed July 29, 2026. Section 4 lists five directions; this record is the second.
\item The supplied summary quotes the authors as having ``failed to reproduce this line of arguments by cropping for three-dimensional bodies'' and reports no verified later analogue. It mentions broad empty-polytope summations without precise references or a summand. An unsuccessful method is not a proof of impossibility, and the summary's tetrahedron language is not imposed on the five-point hull.
\item Provenance flags current primary-literature status and redistribution terms NEEDS\_REVIEW.
\end{itemize}
\end{document}
